% GROK REVIEW::GROK-REV-20261002-1212::2026-10-02 12:12 PT
% Independent comments. Not an LSPV ownership-region rewrite.
\documentclass[11pt]{article}
\usepackage[T1]{fontenc}
\usepackage[utf8]{inputenc}
\usepackage[margin=1in]{geometry}
\usepackage{amsmath,amssymb,booktabs}
\usepackage{hyperref}

\title{Independent review of \texttt{j15marti/pnp}\\[0.3em]
\large GROK-REV-20261002-1212}
\author{Independent reviewer}
\date{2026-10-02 12:12 PT}

\begin{document}
\maketitle

\section*{Status}
P versus NP remains open. This note records comments on the repository at tree
\texttt{cc2d9f62396fa60427f3972af45882f2f8089128}. It does not claim a proof of
\(P\neq NP\), and it does not promote any finite experiment to an asymptotic theorem.

\section*{Corpus reviewed}
Surveyed the repository tree, then read repository-root \texttt{main.tex}
(blob \texttt{fae4d2a04ba6f1d5238ee099146da6b1244e2349}),
\texttt{README.md}, \texttt{AGENTS.md}, \texttt{comments/grok/README.md},
the 2026-10-02 11:26 review, the verified arbiter
\texttt{research/main/2026-10-02\_1158\_MAIN\_1158-R1.tex}
(blob \texttt{187629dc25af58087c7f997eb7c881becadb791e}),
and the prior arbiter \texttt{research/main/2026-10-02\_1129\_MAIN\_1129.tex}.
The newest explorer distill on disk is
\texttt{research/agents/grok/2026-10-02\_1119\_GROK\_H1119\_c1ddc8d1.tex};
no later GROK distill was present. Task files and the STRATFLAT1 case loop were
not re-executed. Ren--Williams, ECCC TR26-118 (12 July 2026), was checked against
the published abstract and against the 11:58 literature audit.

\section*{What is adopted}
\begin{itemize}
\item Global status remains OPEN while any load-bearing arrow is open.
\item MAIN 1158-R1 is the verified arbiter. It supersedes the 11:26 review's citation
of MAIN 1056. MAIN 1129 remains the prior verified snapshot. H1119 remains an explorer
distill, not a second arbiter.
\item LEGALFLAT1 is adopted as a PROVED/EXACT local closure strengthening STRATFLAT1.
If polynomial-round challenge legality may depend on the prior transcript, but legality,
verification, and a polynomial-size strategy evaluator are polynomial time, existence of
a source and winning strategy has the form
\[
\exists R,S\;\forall C\;W(x,R,S,C),\qquad W\in P,
\]
hence lies in \(\Sigma_2^P\). This is a closure, not a separation, and not a theorem
that an arbitrary static SAT decider evaluates \(S\).
\item PBDP1 is adopted only as a hostile representation control. Exponentially expanded
prime products can retain polynomial-size arithmetic straight-line representations.
Explicit expansion is FALSIFIED as a standalone lower-bound mechanism. Source-derived
adaptive bit access remains CANDIDATE and cannot import an unproved BitSLP or PosSLP
separation.
\item The 10:58/11:29 direct-simulation bound remains an upper bound on direct simulation,
not a no-sharing theorem. Ren--Williams TR26-118, randomized parity-decision-tree direct
sums, and QBF strategy-size bounds remain restricted-model results. Transfer class:
useful structural analogy requiring a new source-matched compiler.
\item TRACECLOSE1, PH-PROV-CLOSE1, REACQ1, ACRTREE1, SBACR1, EVACR1, COVERBREAK1,
EXCESS1, STRATFLAT1, and BITRANK1's rejection of complex or real Gram rank remain in force.
\end{itemize}

\section*{Comments, without strengthening}
\begin{enumerate}
\item The 11:26 independent-review section in root \texttt{main.tex} still names MAIN 1056
as the newest arbiter. That sentence is a historical record of that pass. MAIN 1129 and
MAIN 1158-R1 supersede it. This review does not rewrite the historical section.
\item LSPV-4's ownership region already records LEGALFLAT1 as a closure and points to
\texttt{research/agents/lspv-4/2026-10-02\_1146\_LSPV-4\_LEGALFLAT1.tex}. No ownership
region is edited here.
\item The priority-queue item P4 named STRATFLAT1 but not LEGALFLAT1. The living document
now requires a provenance candidate to defeat LEGALFLAT1 as well. That is a label
reconciliation, not a new theorem.
\item PBDP1 does not appear by name in the LSPV-4 ownership region. It is recorded in the
11:58 reconciliation and in this review as a representation-bypass control only.
\item Selecting a PH-PROV-CLOSE1 witness can return a strategy description \(S\). That is
witness selection inside the collapse hypothesis. It is not an arbitrary-algorithm compiler.
\end{enumerate}

\section*{Duplication and contradiction audit}
No contradiction was found between MAIN 1158-R1 and LEGALFLAT1. No contradiction was found
between LEGALFLAT1 and STRATFLAT1: the former adds history-dependent legality and does not
weaken the \(\Sigma_2^P\) flattening. The stronger PROVED/EXACT label applies only inside
the stated polynomial-time legality and evaluator hypotheses. Finite STRATFLAT1 controls
remain COMPUTATIONAL.

\section*{What this review does not do}
It does not edit LSPV-1 through LSPV-4 ownership regions, does not revive retired BOOL
predictors, and does not treat prime provenance, strategy-tree cardinality, adaptive query
depth, witness selection under \(P=NP\), explicit prime-product expansion, or finite
signature diversity as an unrestricted SAT-time measure.

\section*{Single next target}
On one source-legal reconstruction-compatible family, construct \(\Phi\) and the legal
slice matrices and prove or falsify \(\operatorname{rank}_{\mathbb F_2}\Phi=\omega(m)\)
by an intrinsically characteristic-two certificate. On that same family, search for compact
Skolem circuits, algebraic evaluators, decision diagrams, source/history quotients, and
arithmetic/CRT representations. A polynomial representation should be recorded as a hostile
closure of the provenance route, not as progress toward \(P\neq NP\).

% END_OF_VERIFIED_REPORT::GROK::GROK-REV-20261002-1212
\end{document}
